\documentclass[12pt]{article}
\usepackage[utf8]{inputenc}
\usepackage[T1]{fontenc}
\usepackage{hyperref}
\title{Tama Language Resources for K-12 Curriculum Development}
\author{Compiled by Claude Sonnet 5 (claude-sonnet-5) for Andr\'as Kornai}
\date{}
\begin{document}
\maketitle

\section{Introduction}
Tama (ISO 639-3 \texttt{tma}, Glottocode tama1331) is a Taman language spoken in Ouaddai, eastern Chad, and Darfur, Sudan. Sources conflict on speaker numbers (Wikipedia gives 460,000 for 2022--2023; another passage cites 63,000 in Dar Tama). It is described as largely unwritten. Do not confuse it with the Tama of Colombia or other similarly named languages.

\section{Primary Educational Resources}
No primers, textbooks or school programmes were found.

\section{Linguistic and Reference Materials}
\begin{itemize}
\item Tama--English dictionary on Webonary (about 1,422 entries, glosses in English, French and Chadian Arabic): \url{https://www.webonary.org/tama/g11eb0dac-e026-4674-81e8-8c022105d023/}
\item A second Webonary Tama page: \url{https://www.webonary.org/tama/g8b5b5c56-84f1-43d7-a3a6-dca0dea3b8a0/}
\item Wikipedia article on the Tama language: \url{https://en.wikipedia.org/wiki/Tama_language}
\item Wikipedia article on the Taman languages: \url{https://en.wikipedia.org/wiki/Taman_languages}
\item Sudanese doctoral thesis on noun-formation morphology in Tama compared with English, with a proposal for a writing system: \url{https://repository.sustech.edu/handle/123456789/16109}
\item CNRS Reflex bibliography for Tama (lists works by Edgar, Bryan and Bombay): \url{https://reflex.cnrs.fr/Lexiques/webball/biblio.php?LGNOM=Tama}
\end{itemize}

\section{Religious and Community Resources}
None found.

\section{Digital and Audio Resources}
The Webonary dictionary is the only digital resource found.

\section{Curriculum Development Context}
The language appears to have no standard orthography in schools, and no teaching materials turned up. A heritage curriculum would start from the Webonary dictionary and the literature listed in the CNRS bibliography. I did not verify the contents of those works. Note that Tama is described by some sources as a living language with many speakers; the premise of no native speakers should be checked locally.
\end{document}